\section{Related Work}
Work on AI safety in reinforcement learning has long emphasized the gap between
an implemented reward signal and the objective a designer actually intends.
\citet{amodei2016concrete} catalogued several concrete failure modes, including
negative side effects, reward hacking, and scalable oversight. Subsequent simple
environment suites such as AI Safety Gridworlds and related toy domains made
these problems easier to study in controlled settings
\citep{leike2017ai,dennis2019specifying}. Those environments established the
value of small model organisms for safety research, especially when direct
deployment-scale experimentation is infeasible.

More recent work has sharpened the conceptual vocabulary around specification
failure. Analyses of reward hacking and reward tampering distinguish between
proxy exploitation without explicit corruption and cases where the agent can
interfere with the channel that determines reward
\citep{everitt2021reward,pan2022reward,weng2024specification}. That distinction
matters for evaluation: a benchmark that only measures task success can miss
policies that pursue proxy reward through channel manipulation, while a
benchmark that only measures corruption events may miss broader alignment
failures that do not involve direct tampering.

Another relevant line of work concerns the limits of reward learning and online
specification. When reward models or inferred objectives are updated during
training, agents may still exploit misspecification, feedback loops, or weak
oversight \citep{skalse2022pitfalls}. This motivates separating the training
signal from the evaluation objective, even in synthetic tasks, so that
diagnostic results do not collapse into the tautology that agents optimize the
same function used to judge them.

Procedural generation and distribution shift have also become central to RL
benchmark design. Environments that vary layouts, distractors, and local task
structure can distinguish memorization from robust strategy learning
\citep{cobbe2020leveraging}. In the present work, distribution shift is not an
afterthought but an explicit component of the task tuple: a specification
diagnostic benchmark should test not only whether proxy reward is exploitable in
distribution, but also whether such exploitation persists or worsens when the
task configuration changes.

Finally, a growing literature studies external monitors, anomaly detectors, and
other oversight mechanisms for identifying failures from partial behavioral
traces. That literature motivates our monitor-channel formulation, but it also
raises a design hazard: monitor tasks are often accidentally trivialized by
label leakage. We therefore treat the audit view itself as an experimental
object and require explicit leakage checks before reporting monitorability
results. This differs from work that treats monitor features as fixed side
information and focuses primarily on detector architecture.

The closest prior work is thus not a single benchmark or algorithm, but the
combination of safety toy environments, reward tampering analyses, procedural
generalization benchmarks, and oversight research. \framework contributes a
compact synthesis oriented around measurement discipline: a schema that separates
the relevant objects, a metric suite that quantifies several failure modes
together, and a local-first pipeline that ties the paper directly to saved
artifacts.
